% Part: lambda-calculus
% Chapter: representability
% Section: currying

\documentclass[../../../include/open-logic-section]{subfiles}

\begin{document}

\olfileid{lam}{rep}{cur}
\olsection{کَري کول}

يو $\lambd$-تجريد $\lambd[x][M]$ د يوه آرګومېنټ تابع ښيي۔ دا هغه
وخت څرګند محدوديت دے چې څو آرګومېنټه تابع تعريفول غواړو۔ يوه لار دا
ده چې $\lambd$-حساب داسې پراخ کړو چې جوړې، درې‌ګونې او داسې نور پکې
جوړول شوني شي؛ په دې صورت کښې، مثلاً، درې‌ځاييزه تابع
$\lambd[x][M]$ به د خپل آرګومېنټ په توګه درې‌ګونې ته انتظار کوي۔
خو په \emph{کَري کولو} دا کار ډېر آسان دے۔

يوه بېلګه وګورو۔ د يوې شېبې دپاره فرضوو چې په $\lambd$-حساب کښې
$+$ عمل لرو۔ د جمع تابع $2$-ځاييزه ده، يعنې دوه آرګومېنټونه اخلي۔
خو يو $\lambd$-تجريد يوازې يوځاييزې تابعې راکوي: نحو د
$\lambd[(x,y)][(x+y)]$ په څېر عبارتونه نۀ پرېږدي۔ خو هغه يوځاييزه
تابع~$f_x(y)$ څېړلے شو چې $\lambd[y][(x+y)]$ ئې راکوي او په خپل
يو آرګومېنټ~$y$ کښې~$x$ ورزياتوي۔ په حقيقت کښې دا يوه تابع نۀ ده،
بلکې د بېلابېلو ``$x$ ورزيات کړه'' تابعو يوه کورنۍ ده، د هر عدد~$x$
دپاره يوه۔ اوس بله يوځاييزه تابع~$g$ د $\lambd[x][f_x]$ په توګه
تعريفولے شو۔ چې پر آرګومېنټ~$x$ تطبيق شي، $g(x)$ تابع~$f_x$
بېرته ورکوي؛ نو د هغې قيمتونه نورې تابعې دي۔ اوس که $g$ پر $x$ او
بيا پايله پر~$y$ تطبيق کړو، دا تر لاسه کوو:
$(g(x))y = f_x(y) = x+y$۔ په دې ډول يوځاييزه تابع~$g$ د جمع د
دوه‌ځاييزې تابعې هماغه کار کولے شي۔ ``کَري کول'' د دوه‌ځاييزو تابعو
په يوځاييزو تابعو د بدلولو همدغه چل ته وايي، چې قيمتونه ئې پخپله
يوځاييزې تابعې وي۔

اوس د $\lambd$-حساب د نحو دننه يوه سمه بېلګه وګورو۔ تابع
$f(x,y) = x$ څنګه ښيو؟ که داسې تابع تعريفول غواړو چې دوه
آرګومېنټونه اخلي او لومړے بېرته ورکوي،
$\lambd[x][\lambd[y][x]]$ ليکلے شو۔ دا په لفظي توګه هغه تابع ده
چې آرګومېنټ~$x$ اخلي او تابع~$\lambd[y][x]$ بېرته ورکوي۔ تابع
$\lambd[y][x]$ بل آرګومېنټ~$y$ اخلي، خو له پامه ئې غورځوي او تل~$x$
بېرته ورکوي۔ وګورو چې که $\lambd[x][\lambd[y][x]]$ پر دوو
آرګومېنټونو تطبيق کړو، څۀ کېږي:
\begin{align*}
  (\lambd[x][\lambd[y][x]])MN
  \bredone & (\lambd[y][M])N \\
  \bredone & M
\end{align*}

په عمومي ډول، که په يوه ترم~$N$ تعريف شوې تابع د $x_1$، \dots،~$x_n$
پاراميټرونه ولري، داسې ئې ليکلے شو:
$\lambd[x_1][\lambd[x_2][\ldots\lambd[x_n][N]]]$۔ که پرې $n$
آرګومېنټونه تطبيق کړو، نو دا تر لاسه کوو:
\begin{multline*}
  (\lambd[x_1][\lambd[x_2][\ldots\lambd[x_n][N]]]) M_1 \dots M_n \bredone\\
  \begin{aligned}
  \bredone {} & (\Subst{(\lambd[x_2][\ldots\lambd[x_n][N]])}{M_1}{x_1}) M_2
  \dots M_n\\
   \eqs {} & (\lambd[x_2][\ldots\lambd[x_n][\Subst{N}{M_1}{x_1}]]) M_2
            \dots M_n \\
  \vdots & \\
  \bredone {} &\Subst{\Subst{N}{M_1}{x_1}\ldots}{M_n}{x_n}
  \end{aligned}
\end{multline*}
\textbf{د ژباړې د نسخې سمون (OLLAM-001):} سرچينې د وروستۍ کرښې د
تعويض په بدنه کښې P ليکلی و، خو په همدې بيان کښې د تابعې بدنه N
ټاکل شوې ده او P دلته نۀ دے تعريف شوے۔ په وروستۍ کرښه کښې N راوړل
شوے دے؛ نورې نښې او د تعويض ترتيب هماغسې دي۔
وروستۍ کرښه په لفظي توګه دا معنا لري چې د تابعې د تعريف په بدنه کښې
د~$x_i$ پر ځاے~$M_i$ کېږدو؛ دا کټ مټ هغه څه دي چې پر تابعې د څو
آرګومېنټونو د تطبيق پر مهال ئې غواړو۔
\end{document}
